\chapter{Bilan}
\label{chap-bilan}

\section{Les deux versions en chiffres}

\begin{table}[ht]
    \centering
    \begin{tabularx}{\textwidth}{@{}>{\raggedright\arraybackslash}X r r@{}}
        \toprule
                                                 & \textbf{Version 1 (2023)} & \textbf{Version 2 (2026)} \\
        \midrule
        Fichiers source (\fichier{.c})           & 1                         & 14                        \\
        Lignes du plus long fichier              & \num{2666}                & \num{617}                 \\
        Fonctions                                & 36                        & 151                       \\
        Longueur moyenne d'une fonction (lignes) & 67                        & 13                        \\
        Plus longue fonction (lignes)            & 420                       & 77                        \\
        Variables globales modifiables           & 7                         & 0                         \\
        Copies de la logique d'un tour de jeu    & 3                         & 1                         \\
        Fournisseurs de style créés par le code  & 27                        & 1                         \\
        Boucles d'événements (\code{gtk\_main})  & 5 appels                  & 1                         \\
        \addlinespace
        Cas de test automatiques                 & 0                         & 54                        \\
        Parties jouées par la suite de tests     & 0                         & plus de \num{6000}        \\
        Couverture du moteur                     & ---                       & \SI{98.7}{\percent}       \\
        Compilateurs vérifiés en continu         & 0                         & 3                         \\
        Avertissements activés                   & \code{-Wall}              & 14 options, bloquants     \\
        \addlinespace
        Interfaces                               & 1                         & 2                         \\
        Lancement hors du dossier du programme   & non                       & oui                       \\
        Partie rejouable par une graine          & non                       & oui                       \\
        \bottomrule
    \end{tabularx}
    \caption{Comparaison chiffrée des deux versions.}
    \label{tab-comparaison}
\end{table}

Le nombre total de lignes a augmenté : \num{3658} lignes de programme contre \num{2666},
auxquelles s'ajoutent \num{1582} lignes de tests. La refonte n'a pas cherché à écrire moins,
mais à écrire \emph{une seule fois} chaque chose, à un endroit où elle peut être vérifiée.
Le chiffre le plus parlant est sans doute la longueur moyenne d'une fonction, divisée par
cinq : 13~lignes, c'est une fonction que l'on lit d'un regard.

\section{Défauts de la version~1 supprimés}

La relecture complète de la première version, préalable à la refonte, a mis au jour
plusieurs défauts que le jeu masquait à l'usage.

\begin{description}[style=nextline, leftmargin=1.2em]
    \item[Lecture hors limites (niveau difficile)]
          La voisine d'un coup au but, tirée au sort, était lue dans la grille avant le contrôle
          de ses coordonnées (\cref{lst-v1-bug}). Lorsque le coup portait sur un bord, l'indice
          valait $-1$ ou $10$ : comportement indéfini.
    \item[Bateaux oubliés (niveau difficile)]
          Si cette voisine était invalide, la poursuite était abandonnée et le bateau blessé
          oublié. Le niveau \enquote{difficile} l'était donc moins que prévu.
    \item[Boucles d'événements imbriquées]
          Chaque redémarrage appelait \code{gtk\_main()} depuis une fonction de rappel.
    \item[Fenêtre figée]
          Le tour de l'ordinateur s'exécutait sans rendre la main à GTK.
    \item[Chemins figés]
          Images chargées relativement au dossier courant ; sons désignés par un chemin absolu
          propre à une machine.
    \item[Fuites]
          Un fournisseur de style par libellé et une image par mise à jour, dont une partie
          jamais libérée.
    \item[Hasard biaisé]
          Tirages par \code{rand() \% n} (\cref{sec-hasard}).
\end{description}

\begin{lstlisting}[basicstyle=\ttfamily\footnotesize,
                   caption={Version~1 : la case est lue avant d'être contrôlée.},
                   label={lst-v1-bug}]
int random = rand() % 4;

if ((data->player1_grid[coords4[random][0]][coords4[random][1]].is_shot == FALSE) &&
    (coords4[random][0] != -1) &&
    (coords4[random][0] != GRID_SIZE) &&
    (coords4[random][1] != -1) &&
    (coords4[random][1] != GRID_SIZE))
{
    around_hit.x = coords4[random][0];
    around_hit.y = coords4[random][1];
}
else {
    around_hit.hits = 0;   /* la poursuite est abandonnée */
}
\end{lstlisting}

Aucun de ces défauts n'avait été remarqué en jouant. Tous auraient été trouvés par les
outils que la version~2 exécute désormais à chaque modification : le premier par
AddressSanitizer, le deuxième par le test qui exige qu'un bateau touché soit poursuivi.

\section{Ce qui a été conservé}

La refonte ne devait pas changer le jeu. Ont été gardés : les règles et la composition des
flottes ; les trois niveaux et leurs noms ; les écrans de règles et de présentation, avec
leurs textes ; l'identité visuelle --- la marine de Sup Galilée contre les pirates, les
bateaux, les explosions et le missile, redessinés depuis ; les messages de victoire et de
défaite.

Quelques ajouts profitent directement au joueur : les coordonnées autour des grilles, le
score en direct, le dévoilement de la flotte adverse en fin de partie, la possibilité de
rejouer sans relancer le programme, et un ordinateur dont on voit tomber chaque tir.

\section{Limites}

\begin{itemize}
    \item \textbf{L'adversaire difficile n'est pas optimal.} Sa carte de densité compte les
          positions de chaque bateau indépendamment des autres. Un calcul exact énumérerait
          les flottes complètes compatibles avec les observations ; il serait plus coûteux
          pour un gain que le \cref{tab-candidats} laisse supposer faible.
    \item \textbf{Le placement manuel n'existe pas.} Comme dans la version~1, le joueur
          retire sa flotte au hasard jusqu'à satisfaction.
    \item \textbf{Les images et les sons d'origine ont dû être refaits.} Ceux de 2023 avaient
          été réunis sans que leurs sources soient notées, ce qui interdisait de diffuser le
          jeu. La version~2.0 les signalait ; depuis la version~2.1, toutes les images sont
          dessinées et tous les sons synthétisés par un script du dépôt
          (\fichier{tools/generate-assets.sh}), sous la même licence que le code. Les captures
          de l'annexe~\ref{chap-captures} montrent ces nouvelles ressources.
    \item \textbf{Le son n'est vérifié qu'en partie.} En intégration continue, le test de bout
          en bout décode réellement les deux effets sonores, mais vers une sortie factice :
          aucune machine de test n'a de carte son. La restitution elle-même reste à
          contrôler à l'oreille.
    \item \textbf{GTK~3 est en fin de vie.} Le projet reste sur la version enseignée en
          cours ; la vue étant isolée, un portage vers GTK~4 ne toucherait que
          \fichier{view.c}.
\end{itemize}

\section{Perspectives}

L'architecture rend plusieurs prolongements accessibles sans toucher à l'existant.

\begin{description}[style=nextline, leftmargin=1.2em]
    \item[Jeu à deux en réseau]
          C'était l'ambition de la bibliothèque fournie avec le sujet d'origine, organisée autour
          d'un client et d'un serveur. Le modèle actuel, qui exécute un pas par appel et ne
          bloque jamais, s'y prête : un \enquote{joueur distant} n'est qu'une autre source de
          coups, au même titre que l'adversaire artificiel.
    \item[Placement manuel]
          La fonction \code{bs\_board\_place\_ship()} existe et refuse déjà les positions
          illicites ; il reste à lui donner une interface.
    \item[Traduction]
          Tous les textes sont regroupés dans un seul fichier, \fichier{texts.h}. Le passage à
          \emph{gettext} serait mécanique.
    \item[Distribution]
          Un paquet Flatpak rendrait le jeu installable d'un clic sur toute distribution Linux ;
          les métadonnées nécessaires sont déjà dans le dépôt.
\end{description}
